\documentclass[11pt]{article}
\usepackage[margin=1in]{geometry}
\usepackage{amsmath,amssymb,amsthm,hyperref}
\newtheorem{theorem}{Theorem}
\newtheorem{lemma}{Lemma}
\newcommand{\e}{\mathrm e}
\newcommand{\rp}{r_p}
\title{The density of Garcia's quadratic-residue jumps}
\author{Principia Math In Progress}
\date{29 September 2026}
\begin{document}
\maketitle

\paragraph{Task 1 (original problem).}
For a positive integer \(p\), let \(r_p(t)\) denote the least
nonnegative residue of the integer \(t\) modulo \(p\).
Garcia \cite[pp.~5, 22]{Garcia} defines, for \(p=4n-1\),
\[
 \Gamma_n(k)=r_p((k-1)^2)+r_p(k+1-3n),\qquad
 J_n=\#\{k\in\mathbb Z:2\le k\le n+2,\ \Gamma_n(k)\ge p\}.
\]
MathDB \href{https://mathdb.com/p/357172}{357172}, recording
\cite[Conjecture 6.5, p.~29]{Garcia}, asks whether
\(J_n/n\to3/8\) as \(n\to\infty\) through integers for which \(4n-1\)
is prime.

\begin{theorem}
With these definitions, \(J_n/n\to3/8\) as \(n\to\infty\) through
\emph{all} integers. In particular the prime-subsequence conjecture holds.
\end{theorem}

We give all approximation details needed for the proof.
Write \(\e(t)=\exp(2\pi i t)\).

\begin{lemma}\label{distribution}
For integers \(q\to\infty\), the points
\[
 z_{q,j}=\left(\frac jq,\frac{r_q(j^2)}q\right),
 \qquad 0\le j<q,
\]
are uniformly distributed in the unit square: the proportion in any
fixed polygon tends to its area. Including or excluding boundary
segments does not affect this limit.
\end{lemma}
\begin{proof}
For fixed integers \(a,b\), set
\[
 S_q(a,b)=\sum_{j=0}^{q-1}\e\left(\frac{bj^2+aj}{q}\right).
\]
If \(b=0\) and \(a\ne0\), the geometric sum is zero once \(q>|a|\).
If \(b\ne0\), putting \(u=j-k\) modulo \(q\) gives
\begin{align*}
 |S_q(a,b)|^2
 &=\sum_{u\bmod q}\e\left(\frac{bu^2+au}{q}\right)
                    \sum_{k\bmod q}\e\left(\frac{2buk}{q}\right)\\
 &=q\sum_{\substack{u\bmod q\\q\mid2bu}}
                 \e\left(\frac{bu^2+au}{q}\right).
\end{align*}
There are \(\gcd(2b,q)\) terms in the last sum, so
\begin{equation}\label{gauss}
 \frac{|S_q(a,b)|}{q}
 \le\sqrt{\frac{\gcd(2b,q)}q}
 \le\sqrt{\frac{2|b|}{q}}\longrightarrow0.
\end{equation}
Here the gcd uses the absolute value of its first argument.
Thus the average over these points of every nonconstant Fourier
monomial \(\e(ax+by)\) tends to zero, and the constant average is one.

We justify passage from monomials to polygons.
On the circle \(\mathbb R/\mathbb Z\), the Fej\'er kernel
\[
 K_H(t)=\frac1H\left|\sum_{\ell=0}^{H-1}\e(\ell t)\right|^2
\]
is a nonnegative trigonometric polynomial of integral one.
For circle distance \(\|t\|\ge\delta>0\), with \(\delta<1/2\),
\[
 K_H(t)\le\frac1{H\sin^2(\pi\delta)}.
\]
This follows by summing the geometric progression.
If \(f\) is continuous and periodic on the square, its convolution
with \(K_H(x)K_H(y)\) converges uniformly to \(f\).
Indeed uniform continuity controls the integral on
\(\|x\|,\|y\|<\delta\), and the displayed bound makes the remaining
integral tend to zero; the product kernel has integral one.
The convolutions are trigonometric polynomials. The Fourier conclusion
therefore implies
\begin{equation}\label{continuous}
 \frac1q\sum_{j=0}^{q-1}f(z_{q,j})
 \longrightarrow\int_{[0,1]^2}f(x,y)\,dx\,dy.
\end{equation}

A polygon, viewed on the torus, has boundary contained in finitely many
line segments, including any segments on the sides of the square.
The area of an \(\varepsilon\)-neighborhood of this boundary tends
to zero. Continuous lower and upper approximations to its indicator
can be formed by linear ramps in the distance to the polygon and its
complement. They agree with the indicator outside that neighborhood,
take values in \([0,1]\), and have integrals converging to the area.
For a closed polygon \(E\), explicit choices are
\[
 f^-_\varepsilon(z)=\min\{1,d(z,E^c)/\varepsilon\},\qquad
 f^+_\varepsilon(z)=\max\{0,1-d(z,E)/\varepsilon\},
\]
where distance is on the torus. They satisfy
\(f^-_\varepsilon\le\mathbf1_E\le f^+_\varepsilon\).
Apply \eqref{continuous}, then let \(\varepsilon\downarrow0\).
The same boundary-neighborhood bound handles any choice of boundary
inclusion. This proves the lemma.
\end{proof}

\begin{proof}[Proof of the theorem]
Take \(n\ge2\) and set \(j=k-1\). For \(1\le j\le n+1\),
\[
 -p<j+2-3n<0,\qquad
 r_p(j+2-3n)=j+n+1.
\]
Consequently the exact defining inequality for a jump is
\begin{equation}\label{jump}
 r_p(j^2)+j\ge3n-2,\qquad 1\le j\le n+1.
\end{equation}
In the unit square define
\[
 a_p=\frac{n+1}{p}=\frac14+\frac5{4p},\qquad
 b_p=\frac{3n-2}{p}=\frac34-\frac5{4p},
\]
and
\[
 E_p=\{(x,y)\in[0,1]^2:0\le x\le a_p,\ x+y\ge b_p\}.
\]
The point \(z_{p,0}=(0,0)\) is outside \(E_p\), and
\eqref{jump} proves the exact identity
\[
 J_n=\#\{0\le j<p:z_{p,j}\in E_p\}.
\]
This expression includes all endpoints with the same inequality as
the original definition.

To handle the moving region, fix \(0<\varepsilon<1/8\) and put
\begin{align*}
 E&=\{(x,y)\in[0,1]^2:0\le x\le1/4,\ x+y\ge3/4\},\\
 E^\varepsilon
 &=\{(x,y)\in[0,1]^2:0\le x\le1/4+\varepsilon,\
                                  x+y\ge3/4-\varepsilon\}.
\end{align*}
For \(5/(4p)\le\varepsilon\) we have
\(E\subseteq E_p\subseteq E^\varepsilon\).
Both bounding regions are fixed polygons, so Lemma~\ref{distribution}
gives
\[
 \operatorname{area}(E)\le
 \liminf_{n\to\infty}\frac{J_n}{p}
 \le\limsup_{n\to\infty}\frac{J_n}{p}
 \le\operatorname{area}(E^\varepsilon).
\]
Their areas are
\[
 \operatorname{area}(E)=\int_0^{1/4}\left(\frac14+x\right)\,dx
                      =\frac3{32},\qquad
 \operatorname{area}(E^\varepsilon)=\frac32(1/4+\varepsilon)^2.
\]
Letting \(\varepsilon\downarrow0\) proves \(J_n/p\to3/32\).
Since \(p/n=4-1/n\to4\), it follows that \(J_n/n\to3/8\).
\end{proof}

\paragraph{Scope and verification.}
The theorem uses the original definition of jumps and proves the full
conjecture. No unproved source identity, primality, or assumption of
independence between the two coordinates is used. The joint distribution
is proved by the mixed Fourier sums above. Exact finite checks compare
the original count with the transformed inequality and polygon count;
they do not replace the analytic convergence argument.

\paragraph{Task 2 (proposed follow-up).}
Find an explicit power-saving error bound for \(J_n-3n/8\), uniformly
over all integers \(n\).

\paragraph{Task 3 (proposed follow-up).}
Determine how prime powers and square divisors of \(4n-1\) affect the
largest fluctuations of \(J_n-3n/8\).

\paragraph{Task 4 (proposed follow-up).}
Determine limiting jump densities when the interval endpoint and the
linear residue term are replaced by general fixed proportions of \(p\).

\paragraph{Task 5 (proposed follow-up).}
Investigate corresponding joint-distribution and jump-count results
for higher-degree polynomial residues, including which composite
modulus sequences permit such limits.

\begin{thebibliography}{1}
\bibitem{Garcia}
J. Garcia, \emph{Class Number of the Imaginary Quadratic Field and
Quadratic Residues Identities}, arXiv:2301.02951v1 (2023), 31 pages,
\url{https://arxiv.org/abs/2301.02951}.
\end{thebibliography}
\end{document}
